
\subsection*{投稿時第4節}
\Comment{R1.15}
「定義4.2：『productive』と『reachable』が未定義である。『この節は typed local control → finite typed basis の順に進む』という表現もくだけすぎている。」

\Reply
同意する。これらの語は予備事項で定義し，型付き文法は完全な精密化の縮約部分として定義したため，到達可能性・生成可能性の繰り返しの断りは不要となった。引用された非形式的な文も削除した。
\Location{第2節、第5.1節。}

\Comment{R1.16}
「脚注2：『productive』はそれ以前に使われている。脚注ではなく第2節で定義すべきだ。」

\Reply
同意する。到達可能・生成可能を標的文法の精密化より前の第2節で定義し、従来の説明脚注は削除した。
\Location{第2節、第5.1節。}

\Comment{R1.17}
「『到達可能かつ生成可能な型付き非終端記号の集合』と書いてあるが、$\widetilde G$ にそうでない記号があるのか。読者の助けを意図した強調かもしれないが、かえって混乱を招く。」

\Reply
査読者の指摘は正しい。$\widetilde G$ は完全な生成語型付き精密化から縮約部分を取り出したものとして定義されるので、残された記号はすべて到達可能かつ生成可能である。重複する限定を削除した。
\Location{第5.1節。}

\Comment{R1.18}
「冒頭の段落では、異なる出現を異なる型付きコピーとして扱う必要がなぜあるのか、理由が全く示されていない。『何のために、なぜ』を明確にする表現が必要だ。例4.3には『型付けが必要な理由』と付くのに、理由自体が述べられていない。……説明するなら、まず学習アルゴリズムを示し、その完全性証明で『ここでは標的文法が型付けにより精密化されていると仮定し、仮説文法が型付き標的導出を模倣できることを示す』と述べたほうが自然だろう。」

\Reply
提案された順序を採用した。第4節で集合依存のバッチ再構成作用素と保守的逐次学習器を先に定義し，第5節で完全性証明のために生成語型付きの標的文法を導入する。その役割は既定義の仮説文法が型付き標的導出を模倣できることの証明である。外側文脈型付けは削除した（R2.2）。生成語型付けを使う単項規則(U)の箇所では，型付き終端規則 $X\to a$ に対し $h(\omega(X))=h(a)$，型付き二項規則に対し $h(\omega(X))=h(\omega(Y))h(\omega(Z))$ が成り立つ。
\Location{第4節、第5.1節、定理5.3。}

\Comment{R1.19}
「補題4.4：『完全な型付き精密化 $\widetilde G^{\mathrm{full}}$ において』の直後に、$G$ に関する導出の式が現れるのは紛らわしい。不要であるか、文末に移すべきだ。」

\Reply
旧補題4.4は再構成によって削除した。必要な内容は命題5.2に吸収し、言語保存と生成語型の不変量を直接述べ、元の文法と型付き精密化の導出は構文的に区別した。
\Location{命題5.2。}

\Comment{R1.20}
「補題4.8：『$X$ が $\widetilde G$ において到達可能かつ生成可能ならば』という限定も同様に読者を混乱させる。」

\Reply
R1.17の縮約済み精密化の規約によって解決した。旧補題4.8と，到達可能性・生成可能性を重ねて仮定する記述を削除した。
\Location{第5.1節。}

\Comment{R1.21}
「冒頭の段落（$\omega,\chi$ の定義）は、$\mathsf{CS}$ の定義直前に置くとつながりがよいだろう。」
「『特徴標本を次で定義する』というのは、定義2.1の条件を満たすことが直ちに明らかでない限り避けるべきだ。」
「『realizable』をここで定義するのに、それより前から使っており、『realized』は未定義である。」
「『canonical observation word』は $\mathsf{CS}(\widetilde G)$ の定義箇所で定義できるだろう。」

\Reply
同意する。標準生成語 $\omega(X)$ は縮約済み精密化の直後，初出より前に定義し，標準文脈 $\chi(X)$ は有限証拠集合の直前に置いた。$\mathsf{Wit}(\widetilde G)$ は厳密な再構成を証明するまでは証拠集合と呼び，その後に特徴標本と位置づける。未使用の基底，$NT(\cdot)$，$\mathrm{Rule}(\cdot)$，realized/realizable，exhibited などの術語は削除した。
\Location{第5.1節。}

\Comment{R1.22}
「補題4.6は、あまりにも自明な系である。系と呼び直すか削除すべきだ。」
「補題4.7はここまでの議論を理解した読者には自明であり、証明は不要だ。」
「定理4.9が何を示すのかわからない。」
「定理4.12の第一文は削除すべきだ。第二、第三文は未定義の『exhibited』のため理解できない。」

\Reply
旧補題4.6--4.7および定理4.9、4.12は削除した。型付き精密化の議論は、ひとつの正規形定義と、言語保存および生成語型の不変量を述べる命題5.2、それに証拠集合の構成に整理した。
\Location{第5.1節。}

\Comment{R1.23}
「$\alpha$ は記号ではなく文字列である。」
「『とくに $X\to YZ$ ならば』というが、その前件と帰結に関連がない。」
「補題4.5の(iii)について『$G$ の各導出は型注釈を消去する』という文は理解できない。『型注釈を消した $G$ の導出は……を与える』という意味か。」
「$\mathrm{NT}(\widetilde G)$ は9ページの $\widetilde V$ と同じではないか。」
「定義4.11の $\mathrm{Rule}(\widetilde G)$ は9ページの $\widetilde P$ と違うのか。」

\Reply
同意する。不適切な文を削除し，$NT(\widetilde G)$ と $\mathrm{Rule}(\widetilde G)$ はそれぞれ $\widetilde V$ と $\widetilde P$ に統一した。持ち上げの主張は命題5.2にまとめ（R1.19），型付き導出から型注釈を消せば元の文法の導出になると証明中に明記した。
\Location{第5.1節。}

\subsection*{投稿時第5節}
\Comment{R1.24}
「ここで定義された学習アルゴリズムは収束しない。これは論文にとって致命的な誤りである。」

\Reply
同意する。投稿時の学習アルゴリズムは収束せず、これは中心的な正しさの誤りだった。R1.1で説明したように、アルゴリズムを保守的な逐次学習器に置き換え、系5.5で構文的な安定化を証明した。
\Location{第4.2節、系5.5。}

\Comment{R1.25}
「系5.8：この議論は証明として受け入れられない。第5.1節の定義では、学習器の仮説文法は正例を受け取るたびに成長し続ける。」

\Reply
R1.1およびR1.24を参照されたい。証明だけでなく学習器自体を変更した。その結果として系5.5の安定化を得た。
\Location{第4.2節、系5.5。}

\Comment{R1.26}
「注5.1：『規則の健全性』が何を意味するのかわからない。……一文加えるなら『文法は各非終端記号 $[x:u,v]$ から $x$ と代入可能な語の集合が導出されるよう設計する』といった説明が役立つ。定義をもたない形容詞を使うより、数学的性質を簡潔に述べるほうがよい。」
「補題5.3と5.4は定義を繰り返しているだけなので削除すべきだ。」

\Reply
定義の繰り返しである注と補題を削除した。仮説文法では観測された非空因子 $x$ ごとに一つの $\NT{x}$ を置き，出現文脈ごとの記号を用いない（R2.2）。健全性は $\NT{x}\Rightarrow^*w$ ならば $x\equiv_Lw$ かつ $h(w)=h(x)$ という不変量で直接述べる（R1.27）。有限標本の非終端記号が構文合同類全体を生成するとは主張しない。
\Location{第4節（規則(B),(U),(L),(S)と健全性の不変量）。}

\Comment{R1.27}
「補題5.2：第二文の第一の主張は自明で、第二の主張は理解できない。対応する『局所配置』とは何か。『より正確には』に続く部分との関係が不明だ。」

\Reply
旧補題5.2と未定義の「対応する局所配置」は削除した。代わりに定理4.2の証明で
\[
\NT{x}\Rightarrow^*w\quad\Longrightarrow\quad x\equiv_Lw\ \text{かつ}\ h(w)=h(x)
\]
という単純な不変量を用い、導出木の高さに関する帰納法で証明する。単項規則の場合は共通の標本文脈と固定 $h$ の代入可能性を、二分規則の場合は $\equiv_L$ がモノイド合同であることを用いる。開始規則によって仮説文法の健全性が従う。
\Location{定理4.2とその健全性の不変量。}

\Comment{R1.28}
「14ページ1行目：『$G$ が SSBNF なので』は $\widetilde G$ の誤りではないか。」

\Reply
書き直しの際に修正した。完全性証明全体で、対象の文法を明確に指定した。
\Location{定理5.3。}

\Comment{R1.29}
「定理5.7：主張を数学的に書くべきだ。『有限標本が有限型付き再構成基底を露出する』とは何を意味するのか。標本は文字列の集合だが、FTRB は組である。」

\Reply
定理5.4は、必要な仮定を定理文にすべて含めて、
$\mathsf{Wit}(\widetilde G)\subseteq K\subseteq L
\Longrightarrow L(\Bh(K))=L$
という明確な含意を記す。
\Location{定理5.4。}

\subsection*{投稿時第6節}
\Comment{R1.30}
「注6.4の立場を採るなら、文法構成が多項式時間で実行できることはほとんど自明であり、この節の内容は一段落で十分なのではないか。」

\Reply
独立していた計算量の節を第5.3節の短い説明にまとめた。定理5.6では $n_K=\sum_{w\in K}(|w|+1)$ と置き，空語も数える。因子非終端記号は $O(n_K^2)$，二項・単項規則の候補はそれぞれ $O(n_K^3)$ 個であり，明示的構成時間は $O(n_K^4)$ で抑えられる。更新時間については独立した系ではなく本文中に述べる。
\Location{第5.3節。}
